% ===========
%  CHƯƠNG 11
% ===========
\OpenGrid

\ParallelChapter
{Stakeholder Involvement}
{Sự tham gia của các bên liên quan}
{}

\TriPara
  {
    \EN{
      \begin{flushright}
        \emph{Everyone involved must work together to ensure that reasoning and sense making are central foci of high school mathematics programs.}
      \end{flushright}}
  }
  {
    \EN{
      \begin{flushright}
        \emph{Mọi người có liên quan đều phải cùng làm việc để bảo đảm rằng luận và hiểu giữ vị trí trung tâm trong các chương trình toán trung học phổ thông.}
      \end{flushright}}
  }{\NOTE{}  }

\TriPara
  {
    \EN{\noindent This publication sets forth an ambitious vision for the improvement of high school mathematics by refocusing it on reasoning and sense making. This refocusing is not a minor tweaking of the system but a substantial rethinking of the high school mathematics curriculum that requires the engagement of all involved in high school mathematics. Significant effort will be needed to realign the curriculum to focus on reasoning and sense making, to provide teachers with the professional development needed to develop new understanding of the curriculum, and to ensure that all students are given the resources needed to prepare them for our rapidly changing world.}
  }
  {
    \VI{\noindent Ấn phẩm này đưa ra tầm nhìn đầy tham vọng nhằm cải thiện toán trung học phổ thông bằng cách đặt lại trọng tâm vào luận và hiểu. Việc đặt lại trọng tâm này không phải là chỉnh sửa nhỏ đối với hệ thống, mà là sự suy xét lại sâu rộng về chương trình học toán trung học phổ thông, đòi hỏi sự tham gia của mọi người liên quan đến giáo dục toán ở bậc trung học phổ thông. Sẽ cần nỗ lực đáng kể để điều chỉnh lại chương trình học theo hướng tập trung vào luận và hiểu; cung cấp cho giáo viên hoạt động bồi dưỡng chuyên môn cần thiết nhằm phát triển hiểu biết mới về chương trình học; và bảo đảm mọi học sinh đều được cung cấp nguồn lực cần thiết nhằm chuẩn bị cho các em trước thế giới đang thay đổi nhanh chóng.}
  }{\NOTE{}}

\TriPara
  {
    \EN{The following sections outline questions that decision makers can ask themselves or others as they work with different stakeholder groups to increase their engagement in improving high school mathematics. These questions are intended to form the basis for extended reflection on what is currently happening and what needs to happen, and to spur all involved to begin moving toward the goals of this publication.}
  }
  {
    \VI{Các mục sau phác thảo những câu hỏi mà người ra quyết định có thể tự hỏi mình hoặc hỏi người khác khi làm việc với các nhóm bên liên quan khác nhau nhằm tăng cường sự tham gia của họ vào việc cải thiện toán trung học phổ thông. Những câu hỏi này nhằm làm cơ sở cho suy ngẫm sâu rộng về điều hiện đang diễn ra và điều cần diễn ra, đồng thời thúc đẩy mọi người có liên quan bắt đầu tiến tới mục tiêu của ấn phẩm này.}
  }{\NOTE{}}

\ParallelSection
{Students}
{Học sinh}
{}

\TriPara
  {
    \EN{
      \begin{itemize}
        \item Why is the study of mathematics important for high school students?
          \begin{itemize}
            \item Do students see mathematics as a useful, interesting, and creative area of study?
            
            \item Do they consider an understanding of mathematics to be an important tool for their future success?
            
            \item Are they taking as much mathematics as possible so that they will have a firm foundation for their future career or further study?
          \end{itemize}

        \item What do reasoning and sense making mean to high school students?
          \begin{itemize}
            \item Are they actively working to make sense of the mathematics they are learning?
            
            \item Are they developing the reasoning that underlies mathematical procedures?
          \end{itemize}

        \item Do high school students work on their mathematics (in class and at home) thoughtfully?\vspace*{\baselineskip}
          \begin{itemize}
            \item Do they look for relationships and connections, both within mathematics and with other areas of inquiry-including other subjects and cocurricular activities?
            
            \item Do they use mathematical procedures in a mindful way that conveys their understanding of them?
            
            \item Do they work at effectively communicating their mathematical reasoning to their classmates and others?
          \end{itemize}

        \item Do they exhibit a productive attitude toward mathematics and its importance in making decisions about parts of their life outside of school?
      \end{itemize}}
  }
  {
    \VI{
      \begin{itemize}
        \item Vì sao việc học toán quan trọng đối với học sinh trung học phổ thông?
          \begin{itemize}
            \item Học sinh có xem toán học là lĩnh vực học tập hữu ích, thú vị và sáng tạo hay không?
            
            \item Các em có xem hiểu biết toán học là công cụ quan trọng cho thành công tương lai của mình hay không?
            
            \item Các em có học toán nhiều nhất có thể để có nền tảng vững chắc cho nghề nghiệp sau này hoặc cho việc học tập tiếp theo hay không?
          \end{itemize}

        \item Luận và hiểu có ý nghĩa gì đối với học sinh trung học phổ thông?
          \begin{itemize}
            \item Các em có đang chủ động nỗ lực hiểu toán học mà mình đang học hay không?
            
            \item Các em có đang phát triển luận làm nền cho thủ tục toán học hay không?
          \end{itemize}

        \item Học sinh trung học phổ thông có làm việc với toán học một cách có suy xét, cả trên lớp lẫn ở nhà, hay không?
          \begin{itemize}
            \item Các em có tìm kiếm quan hệ và kết nối, cả bên trong toán học lẫn với lĩnh vực tìm hiểu khác---bao gồm môn học khác và hoạt động bổ trợ chương trình học---hay không?
            
            \item Các em có sử dụng thủ tục toán học một cách có ý thức, theo cách thể hiện sự hiểu biết của các em về chúng, hay không?
            
            \item Các em có nỗ lực truyền đạt hiệu quả luận toán học của mình tới bạn cùng lớp và người khác hay không?
          \end{itemize}

        \item Các em có thể hiện thái độ tích cực đối với toán học và đối với tầm quan trọng của toán học trong việc đưa ra quyết định liên quan đến những phương diện đời sống ngoài nhà trường hay không?
      \end{itemize}}
  }
  {
    \NOTE{Nhóm câu hỏi này đặt học sinh vào vị trí trung tâm khi xem xét việc cải thiện toán trung học phổ thông. Chất lượng dạy học không chỉ được nhận ra qua nội dung đã dạy hay điểm số đạt được, mà còn qua cách học sinh nhìn toán, làm việc với toán, giao tiếp về toán và vận dụng toán vào đời sống. Vì vậy, đây không phải là phần phụ thêm ở cuối chương, mà là hệ câu hỏi giúp nhận biết liệu mục tiêu của ấn phẩm có đang đi vào trải nghiệm học tập của học sinh hay không.

    Một câu trả lời tích cực cho câu hỏi vì sao việc học toán quan trọng không chỉ nằm ở việc học sinh nói rằng ``toán quan trọng'', mà ở chỗ các em nhận ra toán giúp mình hiểu thế giới, giải quyết vấn đề và đưa ra quyết định. Các em có thể nêu liên hệ cụ thể giữa toán với học tập, nghề nghiệp hoặc đời sống, đồng thời chủ động học toán để xây dựng nền tảng vững chắc cho tương lai, thay vì chỉ học đủ để vượt qua kì thi. Dấu hiệu cần chú ý là học sinh xem toán như công cụ phát triển lâu dài, chứ không chỉ là môn học bắt buộc.

    Khi hỏi luận và hiểu có ý nghĩa gì đối với học sinh, trọng tâm là liệu các em có tham gia thực chất vào hoạt động toán học hay không. Câu trả lời tích cực thể hiện ở việc học sinh không bằng lòng với đáp số đúng mà muốn hiểu vì sao; các em đặt câu hỏi về ý nghĩa của công thức, phương pháp và kết quả; các em có thể diễn đạt bằng lời cách mình nghĩ khi giải toán. Nếu học sinh chỉ đồng nhất luận với việc viết nhiều bước hơn hoặc trình bày dài hơn, thì điều đang hình thành chưa phải là luận theo tinh thần của sách.

    Câu hỏi về việc học sinh có làm việc với toán một cách có suy xét hay không chuyển trọng tâm sang chất lượng của quá trình học. Câu trả lời tích cực thể hiện ở việc học sinh tìm kiếm liên hệ giữa bài này với bài khác, giữa chủ đề này với chủ đề khác, và giữa toán với lĩnh vực khác; các em biết lựa chọn thủ tục thay vì làm theo máy móc; và khi gặp mâu thuẫn, các em biết xem lại, điều chỉnh hoặc kiểm tra cách làm của mình. Dấu hiệu quan trọng không chỉ là học sinh thực hiện được thủ tục, mà là các em giải thích được vì sao mình dùng thủ tục ấy.

    Câu hỏi về giao tiếp cho thấy luận toán học không phải là hoạt động khép kín trong đầu từng cá nhân. Học sinh phát triển tốt khi các em sẵn sàng trình bày cách nghĩ của mình, có thể diễn đạt bằng lời, hình vẽ, kí hiệu hoặc ví dụ, đồng thời biết lắng nghe và phản hồi luận của bạn học. Nếu lớp học chủ yếu là giáo viên nói còn học sinh chép lại, thì câu trả lời cho nhóm câu hỏi này chưa thể là tích cực.

    Câu hỏi cuối cùng về thái độ cũng cần được hiểu thận trọng. Thái độ tích cực đối với toán không đồng nghĩa đơn giản với việc ``thích toán'', mà là thái độ mang tính xây dựng: sẵn sàng nỗ lực khi gặp khó, không né tránh toán trong lựa chọn học tập hay nghề nghiệp, và nhận ra giá trị của toán trong những quyết định ngoài nhà trường. Dấu hiệu đáng lo là khi học sinh tự tách mình khỏi cơ hội học toán vì cho rằng toán không dành cho mình.

    Có thể xem toàn bộ nhóm câu hỏi này như thước đo sư phạm: học sinh trung học phổ thông đang học toán đúng tinh thần của ấn phẩm khi các em thấy toán có giá trị, chủ động đi vào luận và hiểu, học với sự suy xét, biết truyền đạt luận toán học của mình, và hình thành thái độ đủ vững để tiếp tục dùng toán trong học tập, công việc và đời sống. Theo nghĩa đó, danh sách câu hỏi này không chỉ là công cụ quan sát học sinh, mà còn là cách đánh giá chất lượng thực chất của chương trình toán trung học phổ thông.}
  }

\ParallelSection
  {Families}
  {Gia đình}
  {}

\TriPara
  {
    \EN{
      \begin{itemize}
        \item Why should families consider mathematics important for their high school students?
          \begin{itemize}
            \item Why is a broad preparation in mathematics important?
          \end{itemize}

        \item What is meant by reasoning and sense making?
          \begin{itemize}
            \item Will their students learn the skills they need for future success?
            
            \item Will their students be prepared for success in college?
          \end{itemize}

        \item Do some students possess a ``mathematics gene,'' or can all students be successful in learning to reason mathematically?
        
        \item What courses should families ensure that their students are taking so that their students will receive high-quality preparation in mathematics?
        
        \item What should families look for in their students' classroom?
          \begin{itemize}
            \item What kind of homework should their students be given?
            
            \item What kinds of tests should their students take?
          \end{itemize}

        \item How can families best help their students study mathematics?
          \begin{itemize}
            \item Are families encouraging their students to work for true understanding rather than just for completion of homework assignments?
            
            \item Are families reinforcing the importance of persistence, even in the face of frustration?\vspace{\baselineskip}
          \end{itemize}

        \item What are the most important things that families should do to support the mathematics learning of their high school student?
          \begin{itemize}
            \item Are families encouraging their students to challenge themselves academically?
            
            \item Are families communicating that school should be their students' top priority?
          \end{itemize}\vspace{\baselineskip}

        \item What should family members do if their child shows mathematical talent?
      \end{itemize}}
  }
  {
    \VI{
      \begin{itemize}
        \item Vì sao gia đình nên xem toán là quan trọng đối với con em mình ở trung học phổ thông?
          \begin{itemize}
            \item Vì sao sự chuẩn bị toàn diện về toán học lại quan trọng?
          \end{itemize}

        \item Luận và hiểu có nghĩa là gì?\vspace{\baselineskip}
          \begin{itemize}
            \item Con em họ có học được kĩ năng cần thiết cho thành công sau này hay không?
            
            \item Con em họ có được chuẩn bị để thành công ở bậc đại học hay không?
          \end{itemize}

        \item Có phải chỉ một số học sinh mới có ``gen toán học'', hay mọi học sinh đều có thể thành công trong việc học cách luận toán học?
        
        \item Gia đình cần bảo đảm con em mình học những khóa học nào để các em nhận được sự chuẩn bị chất lượng cao về toán?\vspace{\baselineskip}
        
        \item Gia đình nên chú ý điều gì trong lớp học của con em mình?
          \begin{itemize}
            \item Con em họ nên được giao kiểu bài tập về nhà nào?
            
            \item Con em họ nên làm kiểu bài kiểm tra nào?\vspace{\baselineskip}
          \end{itemize}

        \item Gia đình có thể giúp con em mình học toán tốt nhất bằng cách nào?
          \begin{itemize}
            \item Gia đình có đang khuyến khích con em mình làm việc để đạt sự hiểu thực sự, thay vì chỉ làm xong bài tập về nhà, hay không?
            
            \item Gia đình có đang nhấn mạnh tầm quan trọng của sự bền bỉ, ngay cả khi con em mình nản lòng, hay không?
          \end{itemize}

        \item Những điều quan trọng nhất gia đình nên làm để hỗ trợ việc học toán của con em mình ở trung học phổ thông là gì?
          \begin{itemize}
            \item Gia đình có đang khuyến khích con em mình thử thách bản thân trong học tập hay không?
            
            \item Gia đình có đang truyền đạt rằng việc học ở trường nên là ưu tiên hàng đầu của con em mình hay không?
          \end{itemize}

        \item Thành viên trong gia đình nên làm gì nếu con mình bộc lộ năng khiếu toán học?
      \end{itemize}}
  }
  {
    \NOTE{Nhóm câu hỏi này đặt gia đình vào vị trí lực lượng có ảnh hưởng trực tiếp đến việc học toán của học sinh trung học phổ thông. Việc cải thiện toán học ở bậc học này không chỉ phụ thuộc vào chương trình, giáo viên hay lớp học, mà còn phụ thuộc vào cách gia đình hiểu giá trị của toán, đặt kì vọng học tập, hỗ trợ con em mình khi gặp khó khăn, và mở rộng hay thu hẹp cơ hội học toán của các em.

    Câu hỏi vì sao gia đình nên xem toán là quan trọng không hướng đến câu trả lời chung chung rằng ``toán cần thiết'', mà hướng đến việc nhận ra vai trò lâu dài của toán trong học tập, nghề nghiệp và đời sống. Sự chuẩn bị rộng về toán giúp học sinh không bị khóa sớm vào một lộ trình hẹp, đồng thời giữ mở nhiều cơ hội học tiếp, chọn nghề và tham gia lĩnh vực đòi hỏi năng lực tư duy định lượng, luận chặt chẽ và ra quyết định có cơ sở.

    Khi hỏi luận và hiểu có nghĩa là gì, trọng tâm không nằm ở việc học sinh làm đúng bao nhiêu bài, mà ở việc các em có hiểu ý nghĩa của điều mình làm hay không. Gia đình cần nhận ra rằng thành công lâu dài trong toán không chỉ đến từ việc hoàn thành bài tập, mà từ khả năng giải thích cách nghĩ, đặt câu hỏi, hiểu lí do của thủ tục, và dùng toán để biện giải cho kết quả. Đây cũng là nền tảng chuẩn bị cho thành công sau này và cho việc học ở bậc đại học.

    Câu hỏi về ``gen toán học'' nhằm tháo gỡ niềm tin dễ gây hại: cho rằng chỉ một số học sinh sinh ra đã có khả năng học toán. Theo tinh thần của ấn phẩm, mọi học sinh đều có thể phát triển năng lực luận toán học nếu được đặt trong môi trường học tập có kì vọng phù hợp, nhiệm vụ có chất lượng và sự hỗ trợ đúng lúc. Gia đình vì vậy không nên gán nhãn quá sớm rằng con mình ``có khiếu'' hay ``không có khiếu'' toán, vì những nhãn này có thể ảnh hưởng trực tiếp đến nỗ lực, sự tự tin và cơ hội học tập của học sinh.

    Việc chọn khóa học toán cũng cần được nhìn như quyết định dài hạn. Gia đình không nên chỉ tìm khóa học ít áp lực hơn hay dễ đạt điểm hơn, mà cần xem khóa học đó có duy trì sự chuẩn bị toán học đủ rộng, đủ vững và đủ thách thức hay không. Lộ trình toán tốt phải giúp học sinh tiếp tục phát triển hiểu biết, giữ khả năng học lên, và không tự loại mình khỏi các cơ hội học tập hoặc nghề nghiệp trong tương lai.

    Những điều gia đình quan sát trong lớp học, bài tập về nhà và bài kiểm tra cho thấy chất lượng thực sự của trải nghiệm học toán. Lớp học tốt không chỉ yêu cầu học sinh làm đúng, mà còn tạo cơ hội để các em giải thích, trao đổi và thực hành luận. Bài tập về nhà tốt không chỉ là luyện lại thủ tục, mà còn giúp học sinh kết nối, suy nghĩ và diễn đạt. Bài kiểm tra tốt không chỉ đo đáp số, mà còn cho học sinh cơ hội trình bày cách hiểu và đường luận của mình.

    Gia đình hỗ trợ việc học toán tốt nhất không phải bằng cách làm thay, mà bằng cách nuôi dưỡng cách học đúng. Điều đó bao gồm khuyến khích học sinh hướng tới sự hiểu thực sự, không chỉ hoàn thành bài tập; xem khó khăn và sai lầm như một phần của quá trình học; và củng cố sự bền bỉ khi học sinh nản lòng. Khi gia đình phản ứng với khó khăn bằng sự hỗ trợ có định hướng thay vì phán xét năng lực, học sinh có nhiều cơ hội duy trì niềm tin rằng mình có thể tiếp tục học toán.

    Việc quan trọng nhất gia đình nên làm là giữ kì vọng học tập nghiêm túc, khuyến khích học sinh nhận những thử thách học thuật phù hợp, và truyền đạt rõ rằng việc học ở trường cần được ưu tiên. Sự hỗ trợ này không đồng nghĩa với gây áp lực thành tích, mà là tạo môi trường trong đó nỗ lực học tập được xem là có giá trị, được theo dõi và được nâng đỡ.

    Nếu học sinh bộc lộ năng khiếu toán học, gia đình cần tạo điều kiện để các em được học sâu hơn và gặp cơ hội phù hợp, chẳng hạn câu lạc bộ, dự án, cuộc thi hoặc khóa học nâng cao. Tuy nhiên, việc hỗ trợ năng khiếu không nên biến thành áp lực thành tích quá sớm. Điều cần được bảo vệ là sự phát triển lâu dài của năng lực toán học, cùng với tò mò, niềm vui và sự chủ động trí tuệ của học sinh.

    Nhìn chung, mục ``Gia đình'' cho thấy gia đình không đứng bên ngoài quá trình cải thiện toán trung học phổ thông. Cách gia đình hiểu toán, nói về toán, lựa chọn khóa học, phản ứng trước khó khăn và hỗ trợ cơ hội phát triển đều góp phần quyết định liệu luận và hiểu có trở thành thực chất trong việc học toán của học sinh hay không.}
  }

\ParallelSection
  {Teachers}
  {Giáo viên}
  {}

\TriPara
  {
    \EN{
      \begin{itemize}
        \item Are reasoning and sense making the foci of high school mathematics teachers' instruction every day and for every topic within every course?
        
        \item Do teachers realize and demonstrate the importance of mathematics reasoning habits and content knowledge as life skills that will ensure their students' success in the future, not just as the prerequisite for the next mathematics course that their students may take?\vspace{\baselineskip}
          \begin{itemize}
            \item Do teachers help their students see the wide range of careers that involve mathematics-for example, finance, real estate, marketing, advertising, forensics, and even sports journalism?
            
            \item Do teachers emphasize the practical worth of mathematics for addressing real problems and seek contexts in which mathematics can be seen as a useful and important tool for making decisions?
          \end{itemize}

        \item Do teachers see mathematics as a coherent subject in which the reasons that results are true are as important as the results themselves?
        
        \item Do teachers help students make sense of the mathematics for themselves?
        
        \item Do teachers help students develop a productive disposition toward mathematics?
        
        \item Do teachers strive for balance among the areas of the high school mathematics curriculum outlined in section 2 of this publication and help students find the connections among those areas?
        
        \item Do teachers incorporate technology in ways that enhance students' reasoning and sense making?
        
        \item Do teachers use a range of assessments to monitor and promote reasoning and sense making, both in identifying students' progress and in making instructional decisions?\vspace{\baselineskip}
        
        \item Do teachers recognize the importance of immediately beginning to focus their content and instruction on reasoning and sense making while working with administrators and policy makers toward broadly restructuring the high school mathematics program?\vspace{\baselineskip}
        
        \item Do teachers believe that all students are capable of, and can benefit from a focus on, reasoning and sense making, and are they providing the necessary support for students who need extra support?
          \begin{itemize}
            \item Do teachers ensure that students at all levels have mathematics experiences that are focused on reasoning and sense making?
            
            \item Are teachers advocates for their students, ensuring that schools offer all students the courses that are crucial for their mathematical success?
          \end{itemize}\vspace{\baselineskip}

        \item Do teachers seek out professional development opportunities that will help them gain a better understanding of the reasoning behind mathematical concepts and ways to develop mathematical reasoning in their students?
      \end{itemize}}
  }
  {
    \VI{
      \begin{itemize}
        \item Luận và hiểu có phải là trọng tâm trong hoạt động dạy học của giáo viên toán trung học phổ thông mỗi ngày và với mọi chủ đề trong mọi khóa học hay không?
        
        \item Giáo viên có nhận ra và thể hiện được tầm quan trọng của thói quen luận toán học và tri thức nội dung như kĩ năng sống giúp bảo đảm thành công tương lai cho học sinh, chứ không chỉ như điều kiện tiên quyết cho khóa học toán tiếp theo mà các em có thể theo học, hay không?
        
          \begin{itemize}
            \item Giáo viên có giúp học sinh nhìn thấy phạm vi rộng các nghề nghiệp có liên quan đến toán---chẳng hạn tài chính, bất động sản, tiếp thị, quảng cáo, pháp y, và cả báo chí thể thao---hay không?
            
            \item Giáo viên có nhấn mạnh giá trị thực tiễn của toán trong việc giải quyết vấn đề thực tế và tìm kiếm ngữ cảnh trong đó toán có thể được nhìn như công cụ hữu ích và quan trọng để ra quyết định hay không?
          \end{itemize}

        \item Giáo viên có nhìn toán như môn học nhất quán, trong đó lí do khiến kết quả đúng cũng quan trọng không kém chính kết quả ấy, hay không?\vspace{\baselineskip}
        
        \item Giáo viên có giúp học sinh tự mình hiểu toán hay không?
        
        \item Giáo viên có giúp học sinh hình thành khuynh hướng tích cực đối với toán hay không?
        
        \item Giáo viên có cố gắng giữ sự cân đối giữa các mảng của chương trình học toán trung học phổ thông được trình bày trong Phần 2 của ấn phẩm này, đồng thời giúp học sinh tìm ra mối liên hệ giữa những mảng ấy, hay không?
        
        \item Giáo viên có tích hợp công nghệ theo cách giúp tăng cường luận và hiểu của học sinh hay không?
        
        \item Giáo viên có sử dụng nhiều hình thức đánh giá để theo dõi và thúc đẩy luận và hiểu, cả trong việc nhận diện tiến bộ của học sinh lẫn trong việc đưa ra các quyết định dạy học, hay không?
        
        \item Giáo viên có nhận ra tầm quan trọng của việc ngay lập tức bắt đầu định hướng nội dung và việc giảng dạy của mình vào luận và hiểu, đồng thời làm việc với nhà quản lí và nhà hoạch định chính sách để hướng tới việc tái cấu trúc rộng rãi chương trình toán trung học phổ thông hay không?
        
        \item Giáo viên có tin rằng mọi học sinh đều có khả năng tham gia và đều có thể hưởng lợi từ việc đặt trọng tâm vào luận và hiểu, đồng thời có cung cấp sự hỗ trợ cần thiết cho học sinh cần thêm hỗ trợ, hay không?
          \begin{itemize}
            \item Giáo viên có bảo đảm rằng học sinh ở mọi trình độ đều có trải nghiệm toán học lấy luận và hiểu làm trọng tâm hay không?
            
            \item Giáo viên có là người bảo vệ quyền lợi học tập của học sinh, bảo đảm rằng nhà trường cung cấp cho mọi học sinh khóa học có ý nghĩa thiết yếu đối với thành công trong toán học của các em hay không?
          \end{itemize}

        \item Giáo viên có chủ động tìm kiếm cơ hội bồi dưỡng chuyên môn giúp mình hiểu rõ hơn luận đứng sau khái niệm toán học và cách phát triển luận toán học ở học sinh hay không?
      \end{itemize}}
  }
  {
    \NOTE{Nhóm câu hỏi này đặt giáo viên vào vị trí then chốt trong việc chuyển tầm nhìn của ấn phẩm thành trải nghiệm học tập hằng ngày. Nếu ở mục ``Học sinh'', trọng tâm là những gì có thể quan sát ở người học, và ở mục ``Gia đình'', trọng tâm là điều kiện nâng đỡ từ bên ngoài nhà trường, thì ở đây trọng tâm nằm ở quyết định sư phạm của giáo viên: cách dạy, cách chọn nhiệm vụ, cách đánh giá, cách dùng công nghệ, cách nhìn học sinh và cách tiếp tục học nghề.

    Luận và hiểu không được xem như phần bổ sung cho vài bài học đặc biệt, mà là trọng tâm thường trực trong mọi chủ đề và mọi khóa học. Dấu hiệu tích cực là giáo viên thường xuyên yêu cầu học sinh giải thích vì sao, so sánh cách làm, diễn đạt cách nghĩ và xem xét tính hợp lí của kết quả. Nếu luận và hiểu chỉ xuất hiện ở một số tiết học được chuẩn bị riêng, còn phần lớn thời gian vẫn dành cho làm mẫu và lặp lại, thì tinh thần của ấn phẩm chưa thật sự đi vào việc dạy hằng ngày.

    Việc nhấn mạnh thói quen luận toán học và tri thức nội dung như kĩ năng sống cho thấy toán không nên được trình bày chỉ như điều kiện để đi tiếp sang khóa học sau. Giáo viên cần giúp học sinh thấy toán liên quan đến nhiều nghề nghiệp và nhiều tình huống ra quyết định trong đời sống. Khi toán được đặt trong ngữ cảnh có ý nghĩa, học sinh có thể nhìn toán như công cụ để hiểu, phân tích và hành động trong thế giới, thay vì chỉ như chuỗi yêu cầu học tập cần vượt qua.

    Nhìn toán như môn học nhất quán là nhìn kết quả toán học cùng với lí do khiến chúng đúng. Khi đó, bài học không dừng ở công thức hay kết luận, mà làm lộ ra cấu trúc, giả định và kết nối giữa các ý tưởng. Giúp học sinh tự mình hiểu toán không có nghĩa là để các em tự xoay xở, mà là tạo điều kiện để các em suy nghĩ, thử cách tiếp cận, nhận hỗ trợ đúng lúc và đi tới sự hiểu bằng hoạt động trí tuệ của chính mình.

    Khuynh hướng tích cực đối với toán không đồng nghĩa đơn giản với việc học sinh thích toán. Giáo viên nuôi dưỡng khuynh hướng này khi xem sai lầm như một phần của học tập, khuyến khích thử nghiệm, duy trì kì vọng phù hợp và giúp học sinh bền bỉ khi gặp khó. Thái độ của giáo viên trước lỗi sai, lời giải chưa hoàn chỉnh và tiến độ học khác nhau của học sinh thường tác động trực tiếp đến cách học sinh nhìn năng lực toán học của chính mình.

    Các câu hỏi về cân bằng chương trình và kết nối giữa các mảng nhắc rằng toán trung học phổ thông không nên bị dạy như tập hợp rời rạc của đại số, hình học, hàm số, xác suất hay thống kê. Giáo viên cần giúp học sinh nhìn thấy kết nối giữa những mảng ấy. Chính qua kết nối đó, luận và hiểu mới trở thành năng lực tổ chức tri thức, chứ không chỉ là yêu cầu giải thích trong từng bài riêng lẻ.

    Công nghệ chỉ có giá trị sư phạm khi làm sâu thêm luận và hiểu. Nó không nên được dùng chỉ để tính nhanh hơn hay trình bày đẹp hơn, mà cần giúp học sinh quan sát cấu trúc, thử nghiệm phỏng đoán, kiểm tra kết quả, khảo sát sự biến đổi và so sánh nhiều biểu diễn. Theo nghĩa đó, công nghệ là công cụ của tư duy toán học, không chỉ là công cụ thao tác.

    Đánh giá cũng phải phục vụ luận và hiểu. Giáo viên cần sử dụng nhiều hình thức đánh giá để nhận diện tiến bộ của học sinh, điều chỉnh dạy học và tạo cơ hội cho học sinh giải thích lối luận của mình. Nếu đánh giá chỉ đo đáp số hoặc khả năng thực hiện thủ tục, thì tín hiệu thực sự gửi đến học sinh sẽ trái với mục tiêu đặt luận và hiểu làm trọng tâm.

    Việc tái cấu trúc chương trình toán trung học phổ thông không chỉ là nhiệm vụ ở cấp hệ thống. Giáo viên cần bắt đầu ngay trong nội dung và việc dạy của mình, đồng thời làm việc với nhà quản lí và nhà hoạch định chính sách để mở rộng thay đổi ở cấp chương trình. Như vậy, giáo viên vừa là người thực hành trong lớp học, vừa là tác nhân của cải thiện hệ thống.

    Niềm tin rằng mọi học sinh đều có khả năng luận và hiểu là nền tảng của công bằng trong học toán. Câu trả lời tích cực không nằm ở tuyên bố chung rằng mọi học sinh đều học được, mà ở việc giáo viên giữ kì vọng cao đi kèm hỗ trợ phù hợp. Học sinh ở mọi trình độ cần có cơ hội tham gia trải nghiệm toán học tập trung vào luận và hiểu, và giáo viên cần bảo vệ quyền được tiếp cận khóa học quan trọng cho thành công toán học của các em.

    Câu hỏi về bồi dưỡng chuyên môn nhắc rằng giáo viên không thể phát triển luận toán học ở học sinh nếu chính mình chỉ nắm kĩ thuật bề mặt. Bồi dưỡng chuyên môn cần giúp giáo viên hiểu sâu hơn luận làm nền cho khái niệm toán học, đồng thời học cách thiết kế và dẫn dắt tình huống để học sinh phát triển luận toán học. Đây không chỉ là việc học thêm phương pháp trình bày bài giảng, mà là đào sâu hiểu biết toán học và hiểu biết sư phạm của người dạy.

    Nhìn chung, mục ``Giáo viên'' cho thấy chất lượng chương trình toán trung học phổ thông được quyết định bởi lựa chọn sư phạm lặp lại mỗi ngày. Khi giáo viên giữ luận và hiểu làm trọng tâm, xem toán như lĩnh vực có ý nghĩa, tin vào năng lực của mọi học sinh, dùng công nghệ và đánh giá để nâng đỡ tư duy, đồng thời tiếp tục phát triển chính mình trong tư cách người dạy toán, tầm nhìn của ấn phẩm mới có điều kiện trở thành hiện thực trong lớp học.}
  }

\ParallelSection
  {Administrators}
  {Cán bộ quản lí}
  {}

\TriPara
  {
    \EN{
      \begin{itemize}
        \item Do school districts, schools, departments, and teachers ensure that students receive a high quality high school mathematics curriculum that promotes reasoning and sense making?
        
        \item Do school districts support teachers through long-term professional development that includes reflection on their practice and their work to improve it?
          \begin{itemize}
            \item Do professional development activities help teachers experience mathematics as reasoning and sense making for themselves?
            
            \item Are teachers given time to collaborate during the school day to improve mathematics instruction-with colleagues teaching at the same level to develop tasks that promote reasoning and sense making and with those teaching at higher and lower levels to ensure that reasoning and sense making are being developed across the grades?
            
            \item Do teachers collaboratively analyze students' work and improve the level of formative feedback for students on reasoning and sense making?
            
            \item Is mentoring provided for novice teachers as they work to develop students' reasoning and sense-making skills?
          \end{itemize}

        \item Do schools and school districts work with universities, including mathematicians, statisticians, and mathematics teacher educators, to develop effective programs and courses that address teachers' needs?
        
        \item Do administrators' classroom observations of mathematics teachers focus on reasoning and sense making and provide useful formative feedback?
        
        \item Are support systems in place to provide high school students with the needed assistance to attain high expectations in mathematics?\vspace{\baselineskip}
          \begin{itemize}
            \item Does the school district's counseling infrastructure promote perseverance in taking higher-level mathematics, thereby supporting the goal of increasing students' mathematical reasoning and sense making?
          \end{itemize}
      \end{itemize}}
  }
  {
    \VI{
      \begin{itemize}
        \item Học khu, nhà trường, tổ chuyên môn và giáo viên có bảo đảm rằng học sinh được học theo chương trình toán trung học phổ thông chất lượng cao, thúc đẩy luận và hiểu, hay không?
        
        \item Học khu có hỗ trợ giáo viên thông qua bồi dưỡng chuyên môn dài hạn, trong đó có phản tư về thực hành của mình và nỗ lực cải thiện thực hành ấy, hay không?
          \begin{itemize}
            \item Hoạt động bồi dưỡng chuyên môn có giúp giáo viên tự mình trải nghiệm toán như hoạt động luận và hiểu hay không?
            
            \item Giáo viên có được dành thời gian trong ngày làm việc ở trường để cộng tác nhằm cải thiện việc dạy toán---với đồng nghiệp dạy cùng cấp lớp để phát triển nhiệm vụ thúc đẩy luận và hiểu, và với người dạy ở cấp lớp cao hơn hoặc thấp hơn để bảo đảm rằng luận và hiểu đang được phát triển xuyên suốt các cấp lớp---hay không?
            
            \item Giáo viên có cùng nhau phân tích bài làm của học sinh và nâng cao chất lượng phản hồi quá trình dành cho học sinh về luận và hiểu hay không?
            
            \item Giáo viên mới vào nghề có được cố vấn khi nỗ lực phát triển kĩ năng luận và hiểu của học sinh hay không?
          \end{itemize}

        \item Nhà trường và học khu có phối hợp với trường đại học---bao gồm nhà toán học, nhà thống kê và người đào tạo giáo viên toán---để xây dựng chương trình và khóa học hiệu quả, đáp ứng nhu cầu của giáo viên, hay không?
        
        \item Việc dự giờ giáo viên toán của nhà quản lí có tập trung vào luận và hiểu, đồng thời cung cấp phản hồi quá trình hữu ích, hay không?\vspace{\baselineskip}
        
        \item Có hệ thống hỗ trợ để cung cấp cho học sinh trung học phổ thông sự trợ giúp cần thiết nhằm đạt được kì vọng cao trong toán hay không?
          \begin{itemize}
            \item Hệ thống tư vấn học đường của học khu có khuyến khích học sinh kiên trì theo học toán ở trình độ cao hơn, qua đó hỗ trợ mục tiêu tăng cường luận toán học và hiểu của các em, hay không?
          \end{itemize}
      \end{itemize}}
  }
  {
    \NOTE{Việc phát triển luận và hiểu trong toán trung học phổ thông không thể chỉ được đặt lên vai từng giáo viên. Nó còn phụ thuộc vào điều kiện do học khu, nhà trường, tổ chuyên môn và nhà quản lí tạo ra: chất lượng chương trình học, bồi dưỡng chuyên môn, thời gian cộng tác, cơ chế phản hồi, hỗ trợ cho giáo viên mới, hệ thống tư vấn học đường và quan hệ với các cơ sở đào tạo giáo viên. Vì vậy, nhóm câu hỏi này xem xét liệu hệ thống quản lí có đang tạo điều kiện để luận và hiểu trở thành thực tế trong lớp học hay không.

    Câu hỏi về chương trình toán trung học phổ thông chất lượng cao không chỉ nhắm đến lượng nội dung hay kết quả thi cử. Chương trình có chất lượng, theo tinh thần của ấn phẩm, phải thúc đẩy luận và hiểu: học sinh không chỉ học công thức và thủ tục, mà còn được đặt vào nhiệm vụ đòi hỏi giải thích, kết nối, biện giải và hiểu ý nghĩa toán học. Nếu chương trình chỉ tối ưu tiến độ hoặc điểm số ngắn hạn mà không tạo chiều sâu trong trải nghiệm học toán, thì chưa thể xem là đáp ứng mục tiêu của ấn phẩm.

    Bồi dưỡng chuyên môn dài hạn được nhấn mạnh vì thay đổi bền vững trong dạy học không thể đến từ vài buổi tập huấn rời rạc. Giáo viên cần có cơ hội phản tư về thực hành của mình, thử nghiệm cách làm mới, quan sát kết quả, rồi tiếp tục cải thiện. Quan trọng hơn, bồi dưỡng chuyên môn phải giúp giáo viên tự mình trải nghiệm toán như hoạt động luận và hiểu. Khi giáo viên được đặt vào vị trí người học phải suy nghĩ, giải thích và kết nối, họ mới có cơ sở sư phạm vững hơn để tạo ra trải nghiệm tương tự cho học sinh.

    Câu hỏi về thời gian cộng tác cho thấy luận và hiểu không thể chỉ là nỗ lực cá nhân của từng giáo viên. Giáo viên cần thời gian trong ngày làm việc ở trường để làm việc với đồng nghiệp cùng cấp lớp nhằm thiết kế nhiệm vụ, phân tích bài học và cải thiện việc dạy; đồng thời cần làm việc với người dạy ở cấp lớp cao hơn hoặc thấp hơn để bảo đảm sự phát triển liên tục của luận và hiểu qua các khối lớp. Khi giáo viên cùng phân tích bài làm của học sinh và cải thiện phản hồi quá trình, họ dần hình thành hiểu biết chung về tiến bộ của học sinh trong luận và hiểu.

    Cố vấn cho giáo viên mới cũng là điều kiện quan trọng. Việc cố vấn không chỉ nhằm hỗ trợ quản lí lớp học hay xử lí công việc thường ngày, mà còn giúp giáo viên mới học cách xây dựng lớp học trong đó luận và hiểu giữ vai trò trung tâm. Nếu giáo viên mới chỉ được hướng dẫn cách giữ nề nếp và chạy kịp chương trình, mà không được hỗ trợ để thiết kế nhiệm vụ có chiều sâu và đọc được luận của học sinh, thì mục tiêu cải thiện khó bền vững.

    Sự phối hợp với trường đại học cho thấy nhà trường phổ thông không nên vận hành tách rời khỏi cộng đồng chuyên môn rộng hơn. Làm việc với nhà toán học, nhà thống kê và người đào tạo giáo viên toán có thể giúp học khu và nhà trường xây dựng chương trình, khóa học và hoạt động bồi dưỡng đáp ứng đúng nhu cầu của giáo viên. Sự phối hợp này cũng giúp giữ liên hệ giữa toán phổ thông, toán học chuyên sâu, thống kê và đào tạo giáo viên.

    Câu hỏi về dự giờ của nhà quản lí làm rõ loại tín hiệu mà hệ thống gửi đến giáo viên. Nếu việc dự giờ chỉ chú ý đến trật tự lớp học, tiến độ nội dung hoặc hình thức trình bày, giáo viên sẽ khó xem luận và hiểu là trọng tâm thực sự. Việc quan sát lớp học cần tập trung vào cách học sinh luận, hiểu, giải thích và tham gia vào hoạt động toán học; phản hồi sau dự giờ cũng cần mang tính quá trình, tức giúp giáo viên cải thiện việc dạy thay vì chỉ xếp loại.

    Nhóm câu hỏi cuối về hệ thống hỗ trợ học sinh đưa trọng tâm sang điều kiện công bằng trong học toán. Kì vọng cao chỉ có ý nghĩa khi đi kèm hỗ trợ đủ mạnh để học sinh có thể vươn tới kì vọng ấy. Điều này bao gồm hỗ trợ học tập, tư vấn chọn môn và hệ thống tư vấn học đường khuyến khích học sinh kiên trì theo học toán ở trình độ cao hơn. Nếu hệ thống tư vấn vô tình dẫn học sinh tránh khóa học quan trọng vì sợ khó hoặc sợ thất bại, thì hệ thống ấy đang thu hẹp cơ hội phát triển luận toán học và hiểu của học sinh.

    Nhìn chung, mục này cho thấy chất lượng chương trình toán trung học phổ thông được hình thành không chỉ trong từng lớp học, mà còn qua các quyết định tổ chức của hệ thống. Khi học khu, nhà trường và nhà quản lí tạo điều kiện để giáo viên được bồi dưỡng lâu dài, có thời gian cộng tác, nhận phản hồi đúng trọng tâm, được hỗ trợ khi mới vào nghề, và khi học sinh có hệ thống hỗ trợ để theo đuổi trải nghiệm toán học có chiều sâu, thì luận và hiểu mới có cơ hội trở thành trọng tâm thực sự của chương trình toán trung học phổ thông.}
  }

\ParallelSection
  {Policymakers}
  {Nhà hoạch định chính sách}
  {}

\TriPara
  {
    \EN{
      \begin{itemize}
        \item Why is mathematics important for high school students?
        
        \item Why is high school mathematics education important to the economic competitiveness of the United States in the global marketplace?
        
        \item What is meant by reasoning and sense making in mathematics, and why should these outcomes be crucial foci in high school?
        
        \item Are reasoning and sense making inextricably integrated with content topics in state and district frameworks that guide curricular decisions?\vspace{\baselineskip}
        
        \item Do state and local assessment policies emphasize the need for, and importance of, items that examine students' ability to reason and make sense of mathematical situations?
        
        \item How can policymakers help ensure that all students have access to a rich mathematics curriculum based on reasoning and sense making?
        
        \item Are adequate resources allocated to assist schools and districts in efforts to effectively implement a curriculum based on reasoning and sense making?
      \end{itemize}}
  }
  {
    \VI{
      \begin{itemize}
        \item Vì sao toán học quan trọng đối với học sinh trung học phổ thông?
        
        \item Vì sao giáo dục toán trung học phổ thông quan trọng đối với năng lực cạnh tranh kinh tế của Hoa Kì trên thị trường toàn cầu?
        
        \item Luận và hiểu trong toán có nghĩa là gì, và vì sao hai kết quả này cần là trọng tâm cốt yếu ở trung học phổ thông?
        
        \item Luận và hiểu có được tích hợp không thể tách rời với các chủ đề nội dung trong khung chương trình cấp bang và học khu, tức khung định hướng quyết định chương trình học, hay không?
        
        \item Chính sách đánh giá cấp bang và địa phương có nhấn mạnh nhu cầu và tầm quan trọng của những câu hỏi kiểm tra khả năng học sinh luận và hiểu tình huống toán học hay không?
        
        \item Nhà hoạch định chính sách có thể làm gì để góp phần bảo đảm mọi học sinh đều được tiếp cận chương trình toán phong phú, dựa trên luận và hiểu?
        
        \item Nguồn lực có được phân bổ đầy đủ để hỗ trợ nhà trường và học khu trong nỗ lực triển khai hiệu quả chương trình học dựa trên luận và hiểu hay không?
      \end{itemize}}
  }
  {
    \NOTE{Nhóm câu hỏi dành cho nhà hoạch định chính sách cho thấy cải thiện toán trung học phổ thông không chỉ là vấn đề lớp học hay từng giáo viên. Ở cấp này, câu hỏi không chỉ là học sinh có học tốt hay không, mà là hệ thống giáo dục có tạo đủ điều kiện để mọi học sinh được học toán theo hướng luận và hiểu hay không. Vì vậy, nhóm câu hỏi buộc người làm chính sách nhìn toán như vấn đề chiến lược của giáo dục phổ thông, đồng thời như vấn đề phát triển xã hội và kinh tế.

    Câu hỏi vì sao toán quan trọng đối với học sinh trung học phổ thông không nên được trả lời bằng việc xem toán như môn học để tốt nghiệp hay tuyển sinh. Ở cấp chính sách, toán cần được nhìn như nền tảng cho tư duy phân tích, ra quyết định dựa trên dữ liệu và khả năng thích ứng trước thay đổi nhanh của đời sống hiện đại. Khi chính sách nhìn toán theo hướng này, mục tiêu giáo dục toán vượt khỏi việc bao phủ nội dung hoặc nâng điểm số trước mắt.

    Câu hỏi về năng lực cạnh tranh kinh tế của Hoa Kì mở rộng vấn đề sang bình diện quốc gia. Giáo dục toán trung học phổ thông không chỉ phục vụ từng người học, mà còn liên quan đến chất lượng nguồn nhân lực, năng lực đổi mới công nghệ và khả năng tham gia lĩnh vực có hàm lượng tri thức cao. Theo nghĩa đó, đầu tư cho giáo dục toán là một phần trong chiến lược cạnh tranh dài hạn trên thị trường toàn cầu.

    Khi hỏi luận và hiểu trong toán có nghĩa là gì, và vì sao hai kết quả này cần là trọng tâm cốt yếu ở trung học phổ thông, điều cần nhấn mạnh là luận và hiểu không thể chỉ xuất hiện như khẩu hiệu sư phạm. Chúng phải được xác định như kết quả đầu ra mà hệ thống thật sự muốn học sinh đạt được. Học sinh không chỉ cần làm đúng thủ tục, mà còn cần hiểu, phân tích, kết nối và sử dụng toán trong tình huống mới.

    Câu hỏi về khung chương trình cấp bang và học khu đi thẳng vào vấn đề then chốt: luận và hiểu có được gắn chặt với nội dung hay không. Câu trả lời tích cực là văn bản định hướng không chỉ liệt kê chủ đề phải dạy, mà còn làm rõ rằng nội dung ấy phải được học theo cách nuôi dưỡng luận và hiểu. Nếu điều này không được tích hợp ngay trong khung chương trình, nỗ lực cải thiện ở cấp lớp học sẽ thiếu chỗ dựa học thuật và thể chế, đồng thời dễ bị đẩy lùi bởi áp lực tiến độ hoặc thi cử.

    Câu hỏi về chính sách đánh giá là câu hỏi về tín hiệu hệ thống gửi xuống nhà trường. Nếu chính sách đánh giá coi trọng những câu hỏi đòi hỏi học sinh luận và hiểu tình huống toán học, hệ thống đang nói rõ điều gì là quan trọng. Ngược lại, nếu đánh giá chủ yếu đo khả năng làm thủ tục, mọi tuyên bố về luận và hiểu sẽ khó có sức nặng thực tế. Điều được đánh giá sẽ ảnh hưởng mạnh đến điều được dạy và được học.

    Khi hỏi nhà hoạch định chính sách có thể làm gì để bảo đảm mọi học sinh được tiếp cận chương trình toán phong phú, đoạn này mở vấn đề công bằng ở cấp hệ thống. ``Tiếp cận'' không chỉ là có tên môn học trong danh mục. Nó còn bao gồm việc học sinh có được học với giáo viên đủ năng lực, có lộ trình học hợp lí, và có cơ chế hỗ trợ để không bị loại trừ sớm khỏi cơ hội học toán chất lượng hay không. Chính sách chỉ tuyên bố mở rộng cơ hội nhưng không xử lí điều kiện thực tế thì chưa đủ để tạo ra tiếp cận thực chất.

    Câu hỏi cuối về phân bổ nguồn lực làm rõ rằng mục tiêu chính sách không thể chỉ tồn tại trên giấy. Nguồn lực được phân bổ đầy đủ không chỉ là kinh phí, mà còn là thời gian, học liệu, bồi dưỡng chuyên môn, cơ chế hỗ trợ và điều kiện tổ chức cần thiết để nhà trường và học khu triển khai hiệu quả chương trình học dựa trên luận và hiểu. Nếu không có phân bổ nguồn lực tương xứng, yêu cầu về luận và hiểu rất dễ dừng ở mức tuyên bố.

    Nhìn chung, mục ``Nhà hoạch định chính sách'' khẳng định rằng cải thiện toán trung học phổ thông là quyết định chiến lược cấp hệ thống. Lựa chọn chính sách về khung chương trình, đánh giá, quyền tiếp cận và nguồn lực sẽ quyết định liệu luận và hiểu có trở thành năng lực phổ quát của học sinh hay chỉ tồn tại như mục tiêu trong văn bản. Trách nhiệm của nhà hoạch định chính sách không phải đứng ngoài lớp học, mà là tạo điều kiện để điều tốt nhất có thể diễn ra trong lớp học.}
  }

\ParallelSection
  {Higher Education}
  {Giáo dục đại học}
  {}

\TriPara
  {
    \EN{
      \begin{itemize}
        \item Do colleges follow the recommendations of the \emph{CUPM Curriculum Guide} (CUPM 2004), which suggests a ``move towards reasoning'' as the basis for undergraduate mathematics?
        
        \item Do college entrance or placement examinations test students' ability to reason mathematically, as well as their ability to carry out mathematical procedures?
        
        \item Do mathematics departments act on the recommendations of the \emph{Mathematical Education of Teachers} report (Conference Board of the Mathematical Sciences 2001) for the preparation of high school mathematics teachers?
          \begin{itemize}
            \item In particular, do they provide ``an opportunity for prospective teachers to look deeply at fundamental ideas, to connect topics that often seem unrelated, and to further develop the habits of mind that define mathematical approaches to problems'' (p. 46)?
            
            \item Do they provide courses in statistics and probability for prospective high school teachers that emphasize a data-driven and concept-oriented approach?
            
            \item Do prospective and practicing teachers experience reasoning and sense making in all the major areas of mathematics outlined in the report?
          \end{itemize}

        \item Do colleges of education provide prospective and practicing teachers with courses that emphasize organizing the high school curriculum around reasoning and sense making and instructional methods that support students' development of reasoning and sense making?
        
        \item Do mathematicians, statisticians, and mathematics educators collaborate with one another and with K-12 educators to promote reasoning and sense making in high school mathematics?
          \begin{itemize}
            \item Do they collaborate to develop courses and programs that will support the continued learning of high school teachers as they work to increase the focus in their instruction on reasoning and sense making?
            
            \item Do they seek opportunities to work with K-12 educators in establishing curricular priorities that ensure a balance of content, and in developing state and local curriculum guides and assessment programs that emphasize reasoning and sense making in high school mathematics?
          \end{itemize}
      \end{itemize}}
  }
  {
    \VI{
      \begin{itemize}
        \item Các trường đại học có làm theo khuyến nghị trong \emph{CUPM Curriculum Guide} (CUPM 2004), trong đó đề xuất ``chuyển hướng sang luận'' làm nền tảng cho toán bậc đại học, hay không?
        
        \item Kì thi tuyển sinh hoặc xếp lớp đại học có kiểm tra khả năng luận toán học của học sinh, cũng như khả năng thực hiện thủ tục toán học của các em, hay không?
        
        \item Các khoa toán có thực hiện khuyến nghị trong báo cáo \emph{The Mathematical Education of Teachers} (Conference Board of the Mathematical Sciences 2001) về đào tạo giáo viên toán trung học phổ thông hay không?
          \begin{itemize}
            \item Cụ thể, họ có tạo ``cơ hội để giáo viên tương lai xem xét sâu các ý tưởng nền tảng, kết nối những chủ đề thường có vẻ không liên quan, và tiếp tục phát triển thói quen tư duy định hình cách tiếp cận toán học đối với vấn đề'' (trang 46), hay không?
            
            \item Họ có cung cấp cho giáo viên trung học phổ thông tương lai các khóa học thống kê và xác suất nhấn mạnh cách tiếp cận dựa trên dữ liệu và định hướng khái niệm hay không?
            
            \item Giáo viên tương lai và giáo viên đang hành nghề có được trải nghiệm luận và hiểu trong mọi mảng lớn của toán được phác thảo trong báo cáo hay không?
          \end{itemize}

        \item Các cơ sở đào tạo giáo viên có cung cấp cho giáo viên tương lai và giáo viên đang hành nghề khóa học nhấn mạnh việc tổ chức chương trình trung học phổ thông quanh luận và hiểu, cùng phương pháp dạy học hỗ trợ học sinh phát triển luận và hiểu, hay không?
        
        \item Các nhà toán học, nhà thống kê và nhà giáo dục toán học có cộng tác với nhau và với nhà giáo dục K-12 để thúc đẩy luận và hiểu trong toán trung học phổ thông hay không?\vspace{\baselineskip}
          \begin{itemize}
            \item Họ có cộng tác để phát triển khóa học và chương trình hỗ trợ việc học tập liên tục của giáo viên trung học phổ thông khi giáo viên nỗ lực đặt trọng tâm mạnh hơn vào luận và hiểu trong hoạt động dạy học, hay không?
            
            \item Họ có tìm kiếm cơ hội làm việc với nhà giáo dục K-12 trong việc xác lập ưu tiên chương trình học nhằm bảo đảm cân bằng nội dung, cũng như trong việc phát triển tài liệu hướng dẫn chương trình và chương trình đánh giá cấp bang và địa phương nhấn mạnh luận và hiểu trong toán trung học phổ thông, hay không?
          \end{itemize}
      \end{itemize}}
  }
  {
    \NOTE{Nhóm câu hỏi này cho thấy giáo dục đại học không đứng ngoài nỗ lực cải thiện toán trung học phổ thông. Trường đại học, khoa toán, cơ sở đào tạo giáo viên và cộng đồng học thuật rộng hơn đều góp phần định hình cách giáo viên được chuẩn bị, cách học sinh được đánh giá khi bước vào đại học, và cách hệ thống hiểu vị trí của luận và hiểu trong học toán. Trọng tâm không chỉ là đại học có làm tốt phần việc riêng hay không, mà là đại học có góp phần tạo hệ thống liên kết, trong đó luận và hiểu được coi là trọng tâm từ phổ thông đến sau phổ thông hay không.

    Câu hỏi về \emph{CUPM Curriculum Guide} cho thấy ngay trong toán bậc đại học cũng có yêu cầu chuyển hướng sang luận. Nếu đại học vẫn chủ yếu xem toán như tập hợp kĩ thuật và thủ tục, thì rất khó đòi hỏi phổ thông đi theo hướng khác. Câu trả lời tích cực là trường đại học xem luận như nền tảng của việc học toán đại học, không như phẩm chất phụ thêm dành cho số ít sinh viên giỏi.

    Câu hỏi về kì thi tuyển sinh và xếp lớp nêu một điểm nghẽn hệ thống. Nếu các kì thi này chỉ kiểm tra khả năng làm thủ tục, thông điệp gửi xuống phổ thông sẽ là điều có giá trị vẫn chỉ là thao tác đúng và nhanh. Khi đó, dù phổ thông nói nhiều đến luận và hiểu, trọng tâm ấy vẫn dễ bị suy yếu. Vì vậy, kì thi đại học cần đo được khả năng luận toán học, bên cạnh khả năng thực hiện thủ tục toán học.

    Nhóm câu hỏi về báo cáo \emph{The Mathematical Education of Teachers} đi thẳng vào việc đào tạo giáo viên toán. Giáo viên tương lai không chỉ cần thêm kĩ thuật trình bày bài giảng, mà cần được học toán đủ sâu để thấy ý tưởng nền tảng, nối được chủ đề tưởng như rời nhau, và hình thành thói quen tư duy đặc trưng của toán học. Đào tạo giáo viên toán vì thế không thể chỉ là đào tạo cách dạy, mà còn phải xây dựng hiểu biết toán học đủ mạnh để người dạy có thể nuôi dưỡng luận và hiểu ở học sinh.

    Câu hỏi riêng về thống kê và xác suất làm rõ yêu cầu ấy trong lĩnh vực ngày càng quan trọng. Giáo viên trung học phổ thông tương lai cần được học thống kê và xác suất theo hướng dựa trên dữ liệu và định hướng khái niệm, thay vì chỉ học công thức và kĩ thuật tính. Nếu đại học không chuẩn bị như vậy, việc dạy thống kê và xác suất ở phổ thông khó đạt chiều sâu cần thiết.

    Câu hỏi về cơ sở đào tạo giáo viên đặt ra trách nhiệm sư phạm của giáo dục đại học. Các khóa học cần giúp giáo viên tương lai và giáo viên đang hành nghề hiểu không chỉ dạy nội dung nào, mà còn tổ chức chương trình và chọn phương pháp dạy học ra sao để luận và hiểu trở thành trục của trải nghiệm học toán. Nếu khóa học chỉ tập trung vào quản lí lớp học hoặc mẹo trình bày, trọng tâm của ấn phẩm sẽ khó đi vào thực tế.

    Cụm \emph{collaborate} ở phần cuối cho thấy cải thiện giáo dục toán không thể diễn ra trong các khu vực tách rời: khoa toán theo một hướng, cơ sở đào tạo giáo viên theo một hướng, còn giáo dục K--12 theo hướng khác. Nhà toán học, nhà thống kê, nhà giáo dục toán học và nhà giáo dục K--12 cần cùng xây dựng khóa học, chương trình, ưu tiên nội dung, tài liệu hướng dẫn chương trình và chương trình đánh giá theo cùng định hướng. Nhờ vậy, hệ thống mới duy trì được liên thông giữa nội dung toán học, chuẩn bị giáo viên, dạy học phổ thông và kì vọng đại học.

    Nhìn chung, mục ``Giáo dục đại học'' đặt ra yêu cầu về trách nhiệm liên thông. Nếu đại học không coi luận là trọng tâm, tuyển sinh không đo được luận, đào tạo giáo viên không xây được hiểu biết toán học đủ sâu, và đại học không cộng tác với giáo dục K--12, giáo dục phổ thông sẽ thiếu điểm tựa quan trọng. Ngược lại, khi giáo dục đại học, đào tạo giáo viên và K--12 cùng đồng thuận và cùng hành động, luận và hiểu mới có cơ hội trở thành năng lực xuyên suốt hệ thống giáo dục toán.}
  }

\ParallelSection
  {Curriculum Designers}
  {Nhà thiết kế chương trình học}
  {}

\TriPara
  {
    \EN{
      \begin{itemize}
        \item Do curriculum materials for high school mathematics include a central, regular focus on students' reasoning and sense making that goes beyond the inclusion of isolated supplementary lessons or problems?
        
        \item Does the curriculum, whether integrated or following the course sequence customary in the United States, develop connections among content areas so that students see mathematics as a coherent whole?
        
        \item Is a balance maintained in the areas of mathematics addressed, so that statistics, for example, is more than an isolated unit?
        
        \item Does the curriculum emphasize coherence from one course to the next, demonstrating growth in both mathematical content and reasoning?
        
        \item Does the curriculum incorporate technology in ways that enhance students' mathematical understanding?
        
        \item Is the design of the curriculum materials based on knowledge about how to support the learning of all students?
      \end{itemize}}
  }
  {
    \VI{
      \begin{itemize}
        \item Tài liệu chương trình học cho toán trung học phổ thông có đặt trọng tâm thường trực vào luận và hiểu của học sinh, vượt khỏi việc chỉ đưa vào bài học hoặc bài toán bổ sung rời rạc, hay không?
        
        \item Chương trình học, dù được thiết kế theo mô hình tích hợp hay theo trình tự khóa học quen thuộc ở Hoa Kì, có phát triển kết nối giữa các mảng nội dung để học sinh nhìn toán như chỉnh thể nhất quán hay không?
        
        \item Sự cân bằng giữa các mảng toán được đề cập có được duy trì, để chẳng hạn thống kê không chỉ là đơn vị bài học tách biệt, hay không?
        
        \item Chương trình học có nhấn mạnh tính nhất quán từ khóa học này sang khóa học tiếp theo, qua đó cho thấy sự phát triển cả về nội dung toán lẫn luận, hay không?
        
        \item Chương trình học có đưa công nghệ vào theo cách tăng cường sự hiểu toán học ở học sinh hay không?
        
        \item Việc thiết kế tài liệu chương trình học có dựa trên hiểu biết về cách hỗ trợ việc học của mọi học sinh hay không?
      \end{itemize}}
  }
  {
    \NOTE{Nhóm câu hỏi này cho thấy nhà thiết kế chương trình học giữ vai trò quan trọng trong việc biến luận và hiểu thành khả năng thực tế của lớp học. Tầm nhìn của ấn phẩm không thể chỉ giao cho giáo viên tự bổ sung trong từng bài dạy. Cách tài liệu được thiết kế, nhiệm vụ được chọn, mảng nội dung được nối kết, công nghệ được sử dụng và học liệu được mở cho mọi học sinh đều định hình trước điều có thể xảy ra trong việc học toán.

    Câu hỏi về việc tài liệu chương trình học có đặt luận và hiểu làm trọng tâm thường trực nhằm tránh cách làm quen thuộc: để luận và hiểu xuất hiện như vài bài bổ sung cuối chương hoặc nhiệm vụ nâng cao dành cho số ít học sinh. Câu trả lời tích cực là luận và hiểu được đưa vào cấu trúc thường ngày của bài học và nhiệm vụ, để học sinh thường xuyên phải giải thích, kết nối, suy xét và hiểu vì sao, chứ không chỉ thực hiện đúng bước.

    Câu hỏi về việc chương trình học có giúp học sinh nhìn toán như chỉnh thể nhất quán đi vào yêu cầu cốt lõi của FHSM. Đại số, hình học, hàm số, xác suất và thống kê không nên đứng cạnh nhau như khối rời, mà cần được nối bằng ý tưởng lớn, nhiệm vụ học tập và trình tự bài học. Kết nối ở đây không phải ghi chú bên lề, mà là đặc điểm của chính cấu trúc chương trình học.

    Câu hỏi về cân bằng giữa các mảng nội dung cần được hiểu theo nghĩa mạnh. Chương trình học không đạt cân bằng nếu vài mảng chỉ xuất hiện như đơn vị tách biệt, ít thời lượng và ít chiều sâu. Ví dụ về thống kê nhắc rằng mảng nội dung quan trọng không thể chỉ được thêm vào cho có mặt. Câu trả lời tích cực là mỗi mảng có vị trí tương xứng trong toàn chương trình, cả về thời lượng, độ sâu và vai trò trong việc phát triển luận và hiểu.

    Câu hỏi về tính nhất quán từ khóa học này sang khóa học tiếp theo nhấn mạnh rằng học liệu không nên lặp lại rời rạc hoặc đứt đoạn khi học sinh chuyển khóa học. Chương trình học cần cho thấy sự tiến triển: ý tưởng được mở rộng dần, kết nối được làm sâu thêm, yêu cầu về luận được nâng lên theo thời gian. Nhờ vậy, học sinh thấy mình đang đi trên lộ trình có định hướng, thay vì đi qua từng mảng kiến thức rời rạc.

    Câu hỏi về công nghệ đặt ra chuẩn rõ ràng. Công nghệ chỉ có ý nghĩa khi giúp học sinh hiểu toán sâu hơn: nhìn ra cấu trúc, thử nghiệm phỏng đoán, so sánh biểu diễn, kiểm chứng kết quả hoặc theo dõi sự biến đổi của đối tượng toán học. Nếu công nghệ chỉ dùng để tính nhanh hơn hoặc trình bày đẹp hơn, nó chưa thật sự phục vụ mục tiêu của ấn phẩm.

    Câu hỏi cuối về việc thiết kế tài liệu chương trình học có dựa trên hiểu biết về cách hỗ trợ việc học của mọi học sinh làm rõ chiều kích công bằng trong thiết kế. Tài liệu cần tính đến khác biệt về nền tảng, nhịp độ, cách tiếp cận và mức độ tự tin, đồng thời mở nhiều lối vào để mọi học sinh có thể tham gia luận và hiểu. Nếu học liệu được thiết kế như thể chỉ dành cho một kiểu học sinh, nhiều em sẽ bị đẩy ra ngoài ngay từ đầu.

    Nhìn chung, mục ``Nhà thiết kế chương trình học'' khẳng định rằng luận và hiểu không thể chỉ là điều giáo viên được mong đợi thêm vào sau này. Chúng phải được xây ngay trong cấu trúc học liệu. Khi tài liệu chương trình học đặt luận và hiểu ở vị trí trung tâm, nối kết các mảng nội dung, giữ cân bằng và tính nhất quán, dùng công nghệ như công cụ nhận thức, đồng thời được thiết kế cho việc học của mọi học sinh, lớp học mới có nền tảng vững để đi đúng tinh thần ấn phẩm.}
  }

\ParallelSection
  {Collaboration among Stakeholders}
  {Sự hợp tác giữa các bên liên quan}
  {}

\TriPara
  {
    \EN{
      \begin{itemize}
        \item Are lines of communication among stakeholders being established to promote mathematical reasoning and sense making as central goals of high school mathematics?
        
        \item Is collaboration among stakeholders a central feature of efforts to improve the high school mathematics program?\vspace{\baselineskip}
        
        \item Is the importance of developing ``learning communities'' that collaboratively seek continuing improvement across the educational system recognized?
        
        \item Is each group of stakeholders identifying both long- and short-term action plans in consultation with other stakeholders and beginning to implement those plans?
      \end{itemize}}
  }
  {
    \VI{
      \begin{itemize}
        \item Các kênh giao tiếp giữa các bên liên quan có đang được thiết lập để thúc đẩy luận toán học và hiểu như những mục tiêu trung tâm của toán trung học phổ thông hay không?
        
        \item Sự hợp tác giữa các bên liên quan có phải là đặc điểm trung tâm của nỗ lực cải thiện chương trình toán trung học phổ thông hay không?
        
        \item Tầm quan trọng của việc xây dựng ``cộng đồng học tập'' cùng tìm kiếm cải thiện liên tục trong toàn hệ thống giáo dục có được thừa nhận hay không?
        
        \item Mỗi nhóm bên liên quan có đang xác định cả kế hoạch hành động dài hạn lẫn ngắn hạn thông qua tham vấn với các bên liên quan khác, đồng thời bắt đầu triển khai các kế hoạch ấy, hay không?
      \end{itemize}}
  }
  {
    \NOTE{Nhóm câu hỏi này khép lại Chương 11 bằng cách chuyển trọng tâm từ từng bên liên quan sang quan hệ giữa họ. Vấn đề không chỉ là mỗi nhóm có làm tốt phần việc riêng hay không, mà là toàn hệ thống có tạo được điều kiện để cùng hướng về mục tiêu chung hay không. Theo tinh thần của chương, sự tham gia chỉ thật sự có ý nghĩa khi được chuyển thành kênh trao đổi rõ ràng, hợp tác bền vững, cộng đồng học tập có khả năng tự cải thiện và kế hoạch hành động cụ thể.

    Câu hỏi về kênh trao đổi không chỉ nói đến việc truyền thông tin. Điều quan trọng là có đối thoại hai chiều, thường xuyên và có mục đích giữa giáo viên, nhà quản lí, gia đình, nhà hoạch định chính sách, giáo dục đại học và bên liên quan khác. Khi thiếu kênh trao đổi như vậy, ngay cả thiện chí chung cũng dễ dẫn đến hiểu lệch, ưu tiên rời rạc và hành động thiếu phối hợp.

    Câu hỏi về hợp tác cho thấy cải thiện chương trình toán trung học phổ thông không thể dựa vào nỗ lực đơn lẻ. Hợp tác phải là đặc điểm trung tâm của toàn bộ nỗ lực cải thiện: cùng xác định vấn đề, chia sẻ trách nhiệm, phối hợp hành động và cùng chịu trách nhiệm với mục tiêu đặt luận và hiểu vào trung tâm. Khi đó, các bên liên quan không chỉ cùng tồn tại trong một hệ thống, mà thật sự làm việc như một hệ thống.

    Câu hỏi về ``cộng đồng học tập'' nhấn mạnh nhu cầu hình thành văn hóa học hỏi chung. Cộng đồng học tập không chỉ gặp nhau để xử lí việc trước mắt, mà cùng phản tư, học hỏi và tìm kiếm cải thiện liên tục trên toàn hệ thống giáo dục. Nhờ đó, cải thiện không còn là đợt thay đổi ngắn hạn, mà trở thành quá trình bền vững.

    Câu hỏi cuối đưa vấn đề về thực thi. Hợp tác chỉ có giá trị khi dẫn đến kế hoạch hành động dài hạn và ngắn hạn được xác định rõ, được xây dựng qua tham vấn, và được bắt đầu thực hiện. Mỗi nhóm bên liên quan cần biết mình phải làm gì trước mắt, hướng tới điều gì về lâu dài, và cần phối hợp với ai để đạt mục tiêu chung.

    Nhìn chung, mục này làm rõ điều kiện hệ thống cuối cùng: luận và hiểu khó trở thành trọng tâm thật sự của toán trung học phổ thông nếu các bên liên quan vẫn làm việc như bộ phận rời nhau. Chỉ khi có kênh trao đổi, hợp tác được xem là cốt lõi, cộng đồng học tập được xây dựng, và kế hoạch hành động được cùng nhau thực hiện, mục tiêu của chương mới có thể vượt khỏi văn bản để đi vào thực tiễn của toàn hệ thống.
    }
  }

\ParallelSection
  {Conclusion}
  {Kết luận}
  {}

\TriPara
  {
    \EN{This publication demonstrates that all high school mathematics programs need to be focused on providing all students with the mathematical reasoning and sense-making skills necessary for success in their lives within the context of rich content knowledge. The need to refocus the high school mathematics curriculum has been evident for many years (National Commission on Excellence in Education 1983; Mathematical Sciences Education Board 1989; National Commission on Mathematics and Science Teaching for the 21st Century 2000) and has been reflected in the policies and priorities of the National Council of Teachers of Mathematics for nearly thirty years, from \emph{An Agenda for Action} (1980), to \emph{Curriculum and Evaluation Standards for School Mathematics} (1989), through \emph{Principles and Standards for School Mathematics} (2000a). Although many high school mathematics programs have made significant progress toward achieving such a refocusing, significant work remains to be done to make the deep-rooted, nationwide changes that are vital in meeting the mathematical needs of all students.}
  }
  {
    \VI{Ấn phẩm này cho thấy rằng mọi chương trình toán trung học phổ thông đều cần tập trung vào việc trang bị cho tất cả học sinh các kĩ năng luận toán học và hiểu cần thiết để các em có thể thành công trong cuộc sống, trên nền tảng tri thức nội dung phong phú. Nhu cầu tái định hướng chương trình học toán trung học phổ thông đã trở nên rõ ràng trong nhiều năm qua (National Commission on Excellence in Education 1983; Mathematical Sciences Education Board 1989; National Commission on Mathematics and Science Teaching for the 21st Century 2000) và đã được phản ánh trong chính sách và ưu tiên của Hội đồng Quốc gia Giáo viên Toán học Hoa Kì (NCTM) trong gần ba thập kỉ, từ \emph{An Agenda for Action} (1980), đến \emph{Curriculum and Evaluation Standards for School Mathematics} (1989), rồi qua \emph{Principles and Standards for School Mathematics} (2000a). Mặc dù nhiều chương trình toán trung học phổ thông đã đạt được tiến bộ đáng kể theo hướng tái định hướng này, vẫn còn rất nhiều việc cần làm để tạo ra những thay đổi sâu rộng, bám rễ trên phạm vi toàn quốc, vốn thiết yếu để đáp ứng nhu cầu toán học của tất cả học sinh.}
  }{\NOTE{}}

\TriPara
  {
    \EN{Through this publication, NCTM is providing a framework for thinking about the changes that must be made and for beginning to consider how those changes might be accomplished. However, many issues extending beyond what could be addressed in this publication remain to be answered. Over the coming years, NCTM will continue to provide resources and initiatives that build on this framework. One of NCTM's initial efforts is a set of topic books that set forth additional guidance in particular content areas. Subsequent volumes may address additional issues, such as ensuring equitable experiences for all students in reasoning and sense making. NCTM will also seek to partner with other organizations concerned with high school mathematics.}
  }
  {
    \VI{Thông qua ấn phẩm này, NCTM cung cấp một khung để suy nghĩ về những thay đổi cần thực hiện, cũng như để bắt đầu xem xét những thay đổi ấy có thể được thực hiện như thế nào. Tuy nhiên, vẫn còn nhiều vấn đề vượt ra ngoài phạm vi mà ấn phẩm này có thể xử lí và cần tiếp tục được trả lời. Trong những năm tới, NCTM sẽ tiếp tục cung cấp nguồn lực và sáng kiến được xây dựng dựa trên khung tư duy này. Một trong những nỗ lực ban đầu của NCTM là bộ sách theo chủ đề, đưa ra hướng dẫn bổ sung trong các mảng nội dung cụ thể. Các tập tiếp theo có thể đề cập đến các vấn đề khác, chẳng hạn như việc bảo đảm trải nghiệm công bằng cho tất cả học sinh trong luận và hiểu. NCTM cũng sẽ tìm cách hợp tác với các tổ chức khác có quan tâm đến toán trung học phổ thông.}
  }{\NOTE{}}

\TriPara
  {
    \EN{Although NCTM can take a leadership role, all stakeholders must join forces and work together in meaningful ways to ensure that the continuing story of missed opportunities to significantly improve high school mathematics across the United States will not be told five years from now, let alone in three decades. We simply cannot afford to wait any longer to address the large-scale changes that are needed. The success of our students and of our nation depends on it.}
  }
  {
    \VI{Mặc dù NCTM có thể đảm nhận vai trò lãnh đạo, tất cả các bên liên quan đều phải liên kết lực lượng và hợp tác với nhau một cách thực chất để bảo đảm rằng câu chuyện kéo dài về những cơ hội bị bỏ lỡ trong việc cải thiện đáng kể toán trung học phổ thông trên khắp Hoa Kì sẽ không còn được nhắc lại trong năm năm tới, chứ chưa nói đến ba thập kỉ nữa. Chúng ta không thể tiếp tục chờ đợi thêm nữa để xử lí những thay đổi quy mô lớn đang cần thiết. Thành công của học sinh và của cả quốc gia phụ thuộc vào điều đó.}
  }{\NOTE{}}

\CloseGrid